\section{Path-Space Reduction and Jacobi Response}
\label{sec:path-reduction}

This section separates the algebraic, variational, and differential layers of
restricted geometric control.  The first layer needs only path concatenation.
The second adds compactness and convexity.  The third studies a smooth
nondegenerate optimal path through its second variation.

\subsection{Transition costs and min-plus composition}

Let \(\mathcal M\) be a state space and let
\(\Path_{s,t}(x,y)\) be a declared class of admissible paths from \(x\) at
time \(s\) to \(y\) at time \(t\).  A path has action
\(\mathcal A_{s,t}(\gamma)\in\R\cup\{+\infty\}\).  Empty path classes have
value \(+\infty\).

\begin{assumption}[Composable path system]
\label{ass:path-composition}
For every \(s<r<t\):
\begin{enumerate}[leftmargin=*,itemsep=0.15em]
  \item restricting \(\gamma\in\Path_{s,t}(x,z)\) at time \(r\) produces
  paths in \(\Path_{s,r}(x,\gamma(r))\) and
  \(\Path_{r,t}(\gamma(r),z)\);
  \item any two admissible segments with a common state at time \(r\) can be
  concatenated into an admissible path;
  \item the action is additive under restriction and concatenation.
  \item for every fixed interval \([a,b]\), there is
  \(m_{a,b}>-\infty\) such that
  \(\mathcal A_{a,b}(\gamma)\ge m_{a,b}\) for every admissible segment.
\end{enumerate}
\end{assumption}

Define the transition cost
\begin{equation*}
  c_{s,t}(x,y)
  =
  \inf_{\gamma\in\Path_{s,t}(x,y)}\mathcal A_{s,t}(\gamma).
\end{equation*}

\begin{theorem}[Path-elimination semigroup]
\label{thm:path-dpp}
Under \cref{ass:path-composition}, for every \(s<r<t\),
\begin{equation}
  \boxed{
  c_{s,t}(x,z)
  =
  \inf_{y\in\mathcal M}
  \{c_{s,r}(x,y)+c_{r,t}(y,z)\}.}
  \tag{P1}
\end{equation}
Both sides lie in \(\R\cup\{+\infty\}\), so the min-plus sum is always
defined.  The equality does not require minimizers, smoothness, or convexity.
\end{theorem}

For a terminal cost \(\Gamma:\mathcal M\to\R\cup\{+\infty\}\) bounded
below, define
\[
  V(s,x)
  =
  \inf_y\{c_{s,T}(x,y)+\Gamma(y)\}.
\]

\begin{corollary}[Bellman principle]
\label{cor:path-bellman}
Under the same assumptions,
\begin{equation}
  \boxed{
  V(s,x)
  =
  \inf_y\{c_{s,r}(x,y)+V(r,y)\},
  \qquad s<r<T.}
  \tag{P2}
\end{equation}
Thus Bellman recursion is associativity of path elimination across time
slices.
\end{corollary}

\subsection{The fixed-horizon energy gauge}

Let \(\mathcal C\subset\mathcal M\) be a terminal set and let
\(\Path_T(x,\mathcal C)\) be an admissible class of absolutely continuous
paths on \([0,T]\) starting at \(x\) and ending in \(\mathcal C\).  Suppose
the class is closed under orientation-preserving absolutely continuous
reparameterization and every finite-length path admits its constant-speed
reparameterization within the class.  Define
\[
  D(x,\mathcal C)=\inf_{\gamma\in\Path_T(x,\mathcal C)}\Len(\gamma),
  \qquad
  \mathcal E_T(x,\mathcal C)
  =\inf_{\gamma\in\Path_T(x,\mathcal C)}
  \frac12\int_0^T\norm{\dot\gamma_t}^2\,\dd t.
\]

\begin{theorem}[Length--energy gauge for RDGC]
\label{thm:length-energy-gauge}
Under the preceding reparameterization assumptions,
\begin{equation}
  \boxed{
  \mathcal E_T(x,\mathcal C)
  =\frac{D(x,\mathcal C)^2}{2T}.}
  \tag{P3}
\end{equation}
The identity holds with value \(+\infty\).  When the common value is finite,
a finite-energy path attains the energy minimum if and only if it has constant
speed and attains the length minimum.  Consequently, whenever
the RDGC path class has this reparameterization property,
\[
  \mathcal E_{K,\F,T}(S_0;H)
  =\frac{D_{K,\F}(S_0;H)^2}{2T}.
\]
Thus fixed time is a gauge choice that removes the reparameterization
degeneracy of the length functional without changing its optimal geometric
path.
\end{theorem}

\begin{remark}[A checkable intervention-cube margin]
The convexity hypothesis (MJ2) follows from a uniform Hessian margin.  For
example, suppose \(\mathcal U_0\subset(-\delta,1+\delta)^m\),
\[
  \operatorname{Hess}W_0\succeq\alpha_W g,
  \qquad
  \norm{\operatorname{Hess}w_i}_{\mathrm{op}}\le b_i,
\]
uniformly on \([0,T]\times\mathcal N\), and
\(\alpha_W>(1+\delta)\sum_i b_i\).  Then every \(W_u\) is geodesically
convex on \(\mathcal U_0\).  The analogous condition
\(\alpha_\Phi>(1+\delta)\sum_i c_i\) for
\(\operatorname{Hess}\Phi_0\succeq\alpha_\Phi g\) and
\(\norm{\operatorname{Hess}\phi_i}_{\mathrm{op}}\le c_i\) treats the terminal
potential.  Affine interventions have \(b_i=c_i=0\), as in
\cref{thm:diag-forced-response}.
\end{remark}

\subsection{Existence and Hamilton--Jacobi characterization}

We next give one sufficient analytic realization.  Let
\((\mathcal M,g)\) be a finite-dimensional Hadamard manifold and
\(\F\subset\mathcal M\) a nonempty closed geodesically convex set.  Admissible
paths are absolutely continuous curves in \(\F\) satisfying
\(\dot\gamma(t)\in\mathcal A_t(\gamma(t))\) almost everywhere, and
\[
  \mathcal A_{s,t}(\gamma)
  =
  \int_s^t
  \{L_r(\gamma(r),\dot\gamma(r))+U(r,\gamma(r))\}\,\dd r.
\]

\begin{assumption}[Tonelli--viability conditions]
\label{ass:tonelli}
The admissible velocity multifunction has nonempty closed convex values and a
Borel graph.  In every smooth coordinate chart, the extended integrand
\[
  \overline L(t,G,Z)
  =L_t(G,Z)+U(t,G)+\iota_{\mathcal A_t(G)}(Z)
\]
is a normal integrand: it is Borel measurable in \(t\), lower semicontinuous
in \((G,Z)\) for almost every \(t\), and convex in \(Z\).  For some
\(p>1\), \(a>0\), and \(b\ge0\), uniformly in \((t,G,Z)\),
\[
  \overline L(t,G,Z)\ge a\norm{Z}_G^p-b.
\]
The viability constraint is sequentially closed under uniformly convergent
base paths and weakly convergent chart velocities.  The terminal cost
\(\Gamma\) is proper, lower semicontinuous, and bounded below, and at least one
finite-action competitor exists whenever the value under consideration is
finite.  The path system is composable in the sense of
\cref{ass:path-composition}.
\end{assumption}

\begin{theorem}[Tonelli existence for restricted paths]
\label{thm:path-tonelli}
Under \cref{ass:tonelli}, every finite value \(V(s,x)\) is attained.  If
pointwise geodesic interpolation preserves admissibility and the complete
action plus terminal cost is strictly geodesically convex for paths with a
common initial state, then the minimizing path is unique.
\end{theorem}

For the HJB layer we use a separate, directly checkable control realization.
Suppose \(\F\) is a smooth complete manifold without boundary,
\(\mathcal U\) is a compact metric control space, and admissible paths solve
\[
  \dot G_t=b(t,G_t,a_t)
\]
for measurable controls \(a_t\in\mathcal U\).  Assume \(b\) is continuous,
locally Lipschitz in \(G\) uniformly in \((t,a)\), and tangent to \(\F\).
Let the running cost \(\ell(t,G,a)\) and terminal cost \(\Gamma(G)\) be
continuous and bounded below, and define
\[
  \mathcal H(t,G,p)
  =\sup_{a\in\mathcal U}
  \{\ip{p}{b(t,G,a)}-\ell(t,G,a)\}.
\]

\begin{proposition}[HJB layer with explicit analytic boundary]
\label{prop:path-hjb}
Assume the preceding compact-control realization is nonexplosive on
\([0,T]\), its value \(V\) is finite and continuous, and admissible controls
are closed under concatenation.  Then \(V\) is a viscosity solution of
\begin{equation}
  -\partial_tV(t,G)
  +\mathcal H(t,G,-\dd_GV(t,G))=0,
  \qquad
  V(T,\cdot)=\Gamma.
  \tag{P4}
\end{equation}
At every differentiability point this equation holds classically.  If the
Hamiltonian and terminal data satisfy a comparison principle in a declared
solution class, then \(V\) is the unique viscosity solution in that class.
When the maximizing control is unique, the optimal path satisfies
\[
  \dot G_t=b(t,G_t,a^\star(t,G_t,-\dd_GV(t,G_t)))
\]
at classical points.  Compactness of \(\mathcal U\) gives maximizers and local
velocity compactness; continuity and the uniform local Lipschitz condition
give the short-time realizability used in the viscosity test.  No comparison
or uniqueness claim is made without the stated comparison hypothesis.
\end{proposition}

The layers above have different force.  \Cref{thm:path-dpp} is an exact
algebraic identity, \cref{thm:length-energy-gauge} fixes the time gauge without
changing the optimal geometric path, \cref{thm:path-tonelli} is a direct-method
theorem, and \cref{prop:path-hjb} is an analytic characterization for the
declared compact-control realization, conditional on comparison only for
uniqueness.

\subsection{Jacobi--Schur response}

Let \(\mathcal Z\) be a finite-dimensional visible parameter space and
\(\Hilb_0\) a real Hilbert space of zero-boundary path fluctuations.  A local
trivialization of the admissible path space writes
\[
  \gamma=Rz+\eta,
  \qquad z\in\mathcal Z,
  \quad \eta\in\Hilb_0,
\]
where \(R\) is a fixed continuous right inverse for the visible boundary or
constraint map.  Let
\(\mathfrak A:\mathcal Z\times\Hilb_0\to\R\) be the resulting action.
We use the Riesz isomorphism to identify \(\Hilb_0^*\) with \(\Hilb_0\).

\begin{theorem}[Jacobi--Schur path-response law]
\label{thm:path-schur}
Suppose \(\mathfrak A\in C^3\) near \((z_0,\eta_0)\), where
\(\eta_0\) is the unique local minimizer at \(z_0\).  Assume the vertical
second variation
\[
  \Jac
  =D^2_{\eta\eta}\mathfrak A(z_0,\eta_0)
\]
is represented by a bounded self-adjoint operator satisfying
\[
  \ip{h}{\Jac h}\ge\mu\norm{h}^2,
  \qquad \mu>0,
  \quad h\in\Hilb_0.
\]
Then there are neighborhoods \(Z_0\) of \(z_0\) and \(N_0\) of \(\eta_0\)
such that every \(z\in Z_0\) has a unique minimizing critical point
\(\eta_\star(z)\in N_0\), and this local minimizing section is \(C^2\).  Write
\[
  \Jac(z)=D^2_{\eta\eta}\mathfrak A(z,\eta_\star(z)).
\]
With
\(\Val(z)=\mathfrak A(z,\eta_\star(z))\),
\begin{equation}
  \boxed{
  D\eta_\star
  =
  -\Jac^{-1}\mathfrak A_{\eta z},}
  \tag{P5}
\end{equation}
and
\begin{equation}
  \boxed{
  D^2\Val
  =
  \mathfrak A_{zz}
  -\mathfrak A_{z\eta}\Jac^{-1}\mathfrak A_{\eta z}.}
  \tag{P6}
\end{equation}
All derivatives on the right are evaluated along the minimizing path.  The
effective value curvature is therefore the Schur complement of the vertical
Jacobi block.
\end{theorem}

\begin{corollary}[Coercive kinetic bridge on a Hadamard family]
\label{cor:rdgc-jacobi-bridge}
Let \(\F\) be a finite-dimensional Hadamard manifold realized as a totally
geodesic submanifold of the ambient geometry, and let \(\gamma_0\) be the constant-speed geodesic from
\(x\in\F\) to a fixed endpoint \(q\in\F\).  On zero-boundary fluctuations
\(h\in H_0^1(\gamma_0^*T\F)\), the second variation of kinetic energy is
\[
  I(h,h)=\int_0^T
  \left(\norm{D_th}^2
  -\ip{R(\dot\gamma_0,h)h}{\dot\gamma_0}\right)\dd t
  \ge\int_0^T\norm{D_th}^2\dd t.
\]
Hence the weak Jacobi operator is coercive in the derivative norm on
\(H_0^1\).  If \(q\) is the unique RDGC target projection, then
\cref{thm:length-energy-gauge} identifies \(\gamma_0\) with the energy-gauged
RDGC minimizer.  Every \(C^3\) endpoint perturbation that preserves the fixed
boundary chart therefore satisfies the full Schur formula in
\cref{thm:path-schur}, including its direct visible Hessian term.  Affine
additive-action intensities with fixed endpoint additionally satisfy the
negative-Gram formula in \cref{cor:path-interaction} locally.
\end{corollary}

\subsection{Global Hadamard M\"obius--Jacobi realization}

The abstract Schur law becomes a global theorem on a whole intervention cube
for mechanical actions on Hadamard path space.  Let \((\mathcal N,g)\) be a
finite-dimensional Hadamard manifold, fix \(x_0\in\mathcal N\) and \(T>0\),
and set
\[
  \Omega_{x_0}
  =\{\gamma\in H^1([0,T],\mathcal N):\gamma(0)=x_0\}.
\]
Let the open set \(\mathcal U_0\Subset\mathcal U\subset\R^m\) contain
\([0,1]^m\), and
consider
\[
  W_u=W_0+\sum_{i=1}^m u_iw_i,
  \qquad
  \Phi_u=\Phi_0+\sum_{i=1}^m u_i\phi_i,
\]
\begin{equation}
  \mathcal A_u(\gamma)
  =\int_0^T
  \left(\frac12\norm{\dot\gamma}^2+W_u(t,\gamma(t))\right)\dd t
  +\Phi_u(\gamma(T)).
  \tag{MJ1}
\end{equation}

\begin{theorem}[Global path-space M\"obius--Jacobi theorem]
\label{thm:global-mobius-jacobi}
Fix \(r\ge2\).  Assume the bulk potentials are continuous in \(t\), all bulk
and terminal potentials are \(C^{r+2}\) in the state variable with covariant
derivatives locally uniform in \((t,u)\), and, for every
\(u\in\mathcal U_0\),
\begin{equation}
  \operatorname{Hess}W_u(t,\cdot)\succeq0,
  \qquad
  \operatorname{Hess}\Phi_u\succeq0.
  \tag{MJ2}
\end{equation}
Assume also that the values and gradients of \(W_u(t,\cdot)\) and \(\Phi_u\)
at \(x_0\) are uniformly bounded over \([0,T]\times\mathcal U_0\).
Then the following conclusions hold.

\begin{enumerate}[leftmargin=*,itemsep=0.35em]
  \item For every \(u\in\mathcal U_0\), \(\mathcal A_u\) has a unique global
  minimizer \(\gamma_u\), and \(u\mapsto\gamma_u\) is \(C^r\) on
  \(\mathcal U_0\).  No global lower bound on the potentials is required.
  The minimizer satisfies
  \begin{equation}
    D_t\dot\gamma_u=\operatorname{grad}W_u(t,\gamma_u),
    \qquad
    \dot\gamma_u(T)+\operatorname{grad}\Phi_u(\gamma_u(T))=0.
    \tag{MJ3}
  \end{equation}

  \item On
  \(\mathcal V_u=\{\xi\in H^1(\gamma_u^*T\mathcal N):\xi(0)=0\}\), the
  second variation is
  \begin{equation}
  \begin{aligned}
    \mathfrak B_u(\xi,\eta)
    ={}&\int_0^T\bigl[
    \ip{D_t\xi}{D_t\eta}
    -\ip{R(\xi,\dot\gamma_u)\dot\gamma_u}{\eta}
    +\operatorname{Hess}W_u[\xi,\eta]\bigr]\dd t\\
    &+\operatorname{Hess}\Phi_u[\xi(T),\eta(T)],
  \end{aligned}
  \tag{MJ4}
  \end{equation}
  and it obeys the uniform bound
  \begin{equation}
    \mathfrak B_u(\xi,\xi)
    \ge\int_0^T\norm{D_t\xi}^2\dd t.
    \tag{MJ5}
  \end{equation}
  Hence its weak Jacobi map
  \(\mathfrak J_u:\mathcal V_u\to\mathcal V_u^*\) is a positive
  self-adjoint isomorphism.

  \item Define
  \[
    \Delta_i(\gamma)
    =\int_0^T w_i(t,\gamma(t))\dd t+\phi_i(\gamma(T)),
    \qquad
    \mathfrak f_i^u=D\Delta_i(\gamma_u),
  \]
  and \(\chi_i^u=\mathfrak J_u^{-1}\mathfrak f_i^u\).  Then
  \begin{equation}
    \partial_{u_i}\gamma_u=-\chi_i^u.
    \tag{MJ6}
  \end{equation}
  The response is equivalently the unique solution of
  \begin{equation}
  \begin{aligned}
    J_u\chi_i^u&=\operatorname{grad}w_i(t,\gamma_u),
    &\chi_i^u(0)&=0,\\
    D_t\chi_i^u(T)+\operatorname{Hess}\Phi_u\,\chi_i^u(T)
    &=\operatorname{grad}\phi_i(\gamma_u(T)),
  \end{aligned}
  \tag{MJ7}
  \end{equation}
  where
  \[
    J_u\xi=-D_t^2\xi-R(\xi,\dot\gamma_u)\dot\gamma_u
    +\operatorname{Hess}W_u\,\xi.
  \]
  Thus bulk and terminal intervention forces share one Green--Jacobi inverse.

  \item The value \(\mathcal V(u)=\mathcal A_u(\gamma_u)\) is \(C^r\) and
  \begin{equation}
    \partial_{u_i}\mathcal V=\Delta_i(\gamma_u),
    \qquad
    \partial_{u_i u_j}^2\mathcal V
    =-\mathfrak B_u(\chi_i^u,\chi_j^u).
    \tag{MJ8}
  \end{equation}
  Equivalently,
  \(D_{uu}^2\mathcal V=-\mathfrak F_u^*\mathfrak J_u^{-1}\mathfrak F_u
  \preceq0\), where \(\mathfrak F_ua=\sum_i a_i\mathfrak f_i^u\).
  More precisely,
  \begin{equation}
  \begin{aligned}
    \frac{\norm{\mathfrak F_ua}_{\mathcal V_u^*}^2}
    {\norm{\mathfrak J_u}}
    &\le -a^\top D_{uu}^2\mathcal V(u)a\\
    &\le
    \norm{\mathfrak J_u^{-1}}
    \norm{\mathfrak F_ua}_{\mathcal V_u^*}^2,
  \end{aligned}
  \tag{MJ8a}
  \end{equation}
  and therefore
  \begin{equation}
    \ker(-D_{uu}^2\mathcal V)=\ker\mathfrak F_u,
    \qquad
    \rank(-D_{uu}^2\mathcal V)=\rank\mathfrak F_u.
    \tag{MJ8b}
  \end{equation}

  \item For \(F(A)=\mathcal V(\mathbf1_A)\) and
  \(\zeta_S=\sum_{B\subseteq S}(-1)^{|S|-|B|}F(B)\), every
  \(|S|\le r\) satisfies
  \begin{equation}
    \zeta_S
    =\int_{[0,1]^S}
    \partial_S^{|S|}\mathcal V
    \left(\sum_{i\in S}u_ie_i\right)\prod_{i\in S}\dd u_i.
    \tag{MJ9}
  \end{equation}
  In particular,
  \begin{equation}
  \boxed{
    \zeta_{\{i,j\}}
    =-\int_0^1\!\int_0^1
    \mathfrak B_{se_i+te_j}
    (\chi_i^{se_i+te_j},\chi_j^{se_i+te_j})\,\dd s\,\dd t.}
    \tag{MJ10}
  \end{equation}
  The same formula on the face
  \(u=\mathbf1_C+se_i+te_j\) gives the conditional pair contrast for every
  background \(C\cap\{i,j\}=\varnothing\), with every coordinate outside
  \(C\cup\{i,j\}\) fixed at zero.
\end{enumerate}
\end{theorem}

\begin{corollary}[Arbitrary finite-order Green--Jacobi recursion]
\label{cor:all-order-mobius-jacobi}
Choose a smooth local path chart, let \(I\subseteq[m]\) satisfy
\(1\le|I|\le r\), and write
\(\eta_I^u=\partial_I^{|I|}\gamma_u\), and interpret the displayed higher
derivatives in that chart (equivalently, under one fixed covariant
differentiation convention).  Let \(\Pi(J)\) denote the set of
partitions of \(J\), with \(\Pi(\varnothing)=\{\varnothing\}\), and put
\(\mathscr F(u,\gamma)=D_\gamma\mathcal A_u(\gamma)\).  Then
\begin{equation}
  \eta_I^u=-\mathfrak J_u^{-1}\mathcal R_I^u,
  \tag{MJ11}
\end{equation}
where
\begin{equation}
  \mathcal R_I^u
  =\sum_{\substack{S\subseteq I,\ |S|\le1,\ \pi\in\Pi(I\setminus S)\\
  (S,\pi)\ne(\varnothing,\{I\})}}
  D_u^S D_\gamma^{|\pi|}\mathscr F
  \bigl[\eta_B^u\bigr]_{B\in\pi}.
  \tag{MJ12}
\end{equation}
Thus the same inverse Jacobi operator recursively generates every path
derivative through order \(r\) and, through (MJ9), every Boolean M\"obius
effect of that order.  Since \(r\) may be prescribed arbitrarily, the result
covers any fixed finite order under the corresponding \(C^{r+2}\) regularity;
\(C^\infty\) data give all finite orders simultaneously.  If \(r\ge3\), then
for distinct \(i,j,k\),
\begin{equation}
  \boxed{
  \begin{aligned}
  \partial_{u_i u_j u_k}^3\mathcal V
  ={}&-D_\gamma^3\mathcal A_u[\chi_i,\chi_j,\chi_k]\\
  &+D_\gamma^2\Delta_i[\chi_j,\chi_k]
  +D_\gamma^2\Delta_j[\chi_i,\chi_k]
  +D_\gamma^2\Delta_k[\chi_i,\chi_j].
  \end{aligned}}
  \tag{MJ13}
\end{equation}
\end{corollary}

The intermediate chart derivatives \(\eta_I^u\) and the displayed recursive
representation depend on the chosen path chart, or equivalently on the fixed
covariant-differentiation convention.  The scalar derivatives
\(\partial_I^{|I|}\mathcal V\) and the resulting Boolean M\"obius effects are
intrinsic.

\subsection{Hard-target response on a smooth active stratum}

We now differentiate the terminal constraint itself.  This gives the missing
hard-target counterpart of the Robin response above.  Let
\(c:\mathcal U\times\mathcal N\to\R\) be smooth and consider
\begin{equation}
  \mathcal V_{\rm hard}(u)
  =\inf\{\mathcal A_u(\gamma):
  \gamma\in\Omega_{x_0},\ c_u(\gamma(T))\le0\}.
  \tag{HT1}
\end{equation}
The statement is local in the parameter and in one smooth active stratum;
the inequality itself remains hard.

\begin{theorem}[Bordered Jacobi response for an active hard target]
\label{thm:hard-target-jacobi}
Fix \(u_0\in\mathcal U\).  Suppose that \(\mathcal A_u\) and \(c_u\) are
\(C^{r+2}\), \(r\ge2\), in compatible local path and state charts.  Let
\(\gamma_0\) be a constrained local minimizer of \textup{(HT1)} with active
constraint, and suppose that there is a strictly positive multiplier
\(\lambda_0\) such that, for
\[
  \mathscr L(u,\gamma,\lambda)
  =\mathcal A_u(\gamma)+\lambda c_u(\gamma(T)),
\]
the KKT equations are
\begin{equation}
  D_\gamma\mathscr L(u_0,\gamma_0,\lambda_0)=0,
  \qquad c_{u_0}(\gamma_0(T))=0.
  \tag{HT2}
\end{equation}
On
\(\mathcal V_0=\{\xi\in H^1(\gamma_0^*T\mathcal N):\xi(0)=0\}\), set
\[
  \mathfrak J_0=D_{\gamma\gamma}^2\mathscr L,
  \qquad
  \mathfrak B_0\xi
  =\dd c_{u_0}(\gamma_0(T))[\xi(T)].
\]
Assume that \(\mathfrak J_0:\mathcal V_0\to\mathcal V_0^*\) is bounded,
self-adjoint, and coercive, and that \(\mathfrak B_0\ne0\).  Then:
\begin{enumerate}[leftmargin=*,itemsep=0.3em]
  \item there is a neighborhood \(U_0\) of \(u_0\) and a unique \(C^r\)
  branch \(u\mapsto(\gamma_u,\lambda_u)\) of active KKT points near
  \((\gamma_0,\lambda_0)\), with \(\lambda_u>0\); every \(\gamma_u\) is the
  unique constrained local minimizer in a common path neighborhood;

  \item on \(\mathcal Y_u=\mathcal V_u\times\R\), the bordered Jacobi--KKT
  map
  \begin{equation}
    \mathfrak K_u
    =
    \begin{pmatrix}
      \mathfrak J_u&\mathfrak B_u^*\\
      \mathfrak B_u&0
    \end{pmatrix}
    :\mathcal Y_u\longrightarrow\mathcal Y_u^*
    \tag{HT3}
  \end{equation}
  is a self-adjoint isomorphism, where
  \(\mathfrak J_u=D_{\gamma\gamma}^2\mathscr L\) and
  \(\mathfrak B_u\xi=\dd c_u(\gamma_u(T))[\xi(T)]\);

  \item if
  \[
    \mathfrak b_i^u
    =D_{(\gamma,\lambda)}\partial_{u_i}\mathscr L
      (u,\gamma_u,\lambda_u)\in\mathcal Y_u^*,
  \]
  then the path and multiplier response and the constrained value Hessian are
  \begin{equation}
    \boxed{
    \partial_{u_i}(\gamma_u,\lambda_u)
    =-\mathfrak K_u^{-1}\mathfrak b_i^u,}
    \tag{HT4}
  \end{equation}
  \begin{equation}
    \boxed{
    \partial_{u_i u_j}^2\mathcal V_{\rm hard}
    =\partial_{u_i u_j}^2\mathscr L
    -\left\langle
      \mathfrak b_i^u,\mathfrak K_u^{-1}\mathfrak b_j^u
    \right\rangle_{\mathcal Y_u^*,\mathcal Y_u}.}
    \tag{HT5}
  \end{equation}
  The direct derivative in \textup{(HT5)} holds \((\gamma,\lambda)\) fixed.
\end{enumerate}
If \(\mathcal A_u\) and \(c_u\) are affine in \(u\), the direct term in
\textup{(HT5)} vanishes.  If the same active branch and bordered
invertibility persist on a Boolean two-face, its exact conditional pair
effect is
\begin{equation}
  \boxed{
  \zeta_{\{i,j\}\mid C}
  =-\int_0^1\!\int_0^1
  \left\langle
  \mathfrak b_i^u,\mathfrak K_u^{-1}\mathfrak b_j^u
  \right\rangle\,\dd s\,\dd t,\quad
  u=\mathbf1_C+se_i+te_j.}
  \tag{HT6}
\end{equation}
Here every coordinate outside \(C\cup\{i,j\}\) is fixed at zero.
Unlike the unconstrained Gram formula, the bordered inverse is indefinite, so
hard-target interactions need not have one sign.
\end{theorem}

The terminal contribution to \(\mathfrak J_u\) is
\(\lambda_u\operatorname{Hess}c_u\), while
\(\mathfrak B_u\) linearizes motion of the active boundary.  Thus
\textup{(HT3)} couples the interior Jacobi equation, endpoint transversality,
and active-target motion in one invertible system.

For fixed endpoints, the same theorem holds on Dirichlet fluctuations and the
terminal forces disappear.  For a general regular Lagrangian it remains valid
when existence, uniqueness, smoothness, and uniform coercivity of the complete
second variation are verified.  At active-set changes the correct object is a
normal-cone graphical derivative coupled to a Jacobi variational inequality;
a classical Hessian need not exist.

Accordingly, \cref{thm:global-mobius-jacobi} is the smooth mechanical
extension of the RDGC path framework, while
\cref{thm:hard-target-jacobi,cor:hard-condition-response} gives the classical
hard-target response on a fixed smooth active stratum.  Neither theorem
asserts classical differentiability through eigenvalue multiplicities or
active-index changes.

The theorem applies to a smooth integral action when its second-variation
bilinear form is continuous and coercive on the chosen Sobolev fluctuation
space.  The Riesz map identifies that bilinear form with \(\Jac\); in a
classical Euler--Lagrange problem this is the weak Jacobi operator with the
declared boundary conditions.

\begin{corollary}[Associative second-order path elimination]
\label{cor:path-schur-associativity}
If \(\eta=(\eta_1,\eta_2)\) and the complete vertical Hessian is coercive,
then eliminating \(\eta_2\) and subsequently \(\eta_1\) gives the same value
Hessian as eliminating \((\eta_1,\eta_2)\) in one step:
\begin{equation}
  \operatorname{Schur}_{(\eta_1,\eta_2)}\mathbb H
  =
  \operatorname{Schur}_{\eta_1}
  \bigl(\operatorname{Schur}_{\eta_2}\mathbb H\bigr).
  \tag{P7}
\end{equation}
\end{corollary}

\subsection{Interaction curvature generated by path relaxation}

Let \(u\in U\subset\R^m\) parameterize smooth mechanism intensities and
suppose the unreduced action is affine in \(u\):
\[
  \mathfrak A(\eta,u)
  =
  \mathfrak A_0(\eta)+\sum_{i=1}^m u_i\Delta_i(\eta).
\]
Define \(\mathcal G:\R^m\to\Hilb_0^*\) by
\[
  \mathcal G a=\sum_i a_iD_\eta\Delta_i.
\]

The finite-dimensional form of the following Schur identity is established
in the companion hidden-state reduction theory
\cite{li2026optimizationgeometrodynamics}.  We record its path-Hilbert form to
connect the abstract reduction formula to the Jacobi and Green operators
constructed above; the generic negative-Gram algebra is used as an established
ingredient.

\begin{corollary}[Path interaction curvature]
\label{cor:path-interaction}
Under \cref{thm:path-schur}, the local reduced value satisfies
\begin{equation}
  \boxed{
  D^2_{uu}\Val
  =
  -\mathcal G^*\Jac^{-1}\mathcal G
  \preceq0.}
  \tag{P8}
\end{equation}
In components,
\begin{equation}
  \partial^2_{u_i u_j}\Val
  =
  -\ip{D_\eta\Delta_i}
  {\Jac^{-1}D_\eta\Delta_j}.
  \tag{P9}
\end{equation}
If the same minimizing branch exists smoothly and the vertical second
variation remains uniformly coercive on an open set containing
\([0,1]^m\), then for two Boolean endpoint settings, with all other
intensities fixed at zero, the mixed finite difference is exactly
\begin{equation}
\begin{aligned}
  &\Val(e_i+e_j)-\Val(e_i)-\Val(e_j)+\Val(0)\\
  &\quad=
  -\int_0^1\!\int_0^1
  \ip{D_\eta\Delta_i}
  {\Jac^{-1}D_\eta\Delta_j}_{\eta_\star(se_i+te_j)}
  \,\dd s\,\dd t.
\end{aligned}
\tag{P10}
\end{equation}
Thus hidden path relaxation creates an observable interaction whenever the two
path perturbations fail to be orthogonal in the inverse-Jacobi metric.  A zero
integrated interaction can also arise through cancellation; only pointwise
zero curvature implies pointwise orthogonality.
\end{corollary}

